\subsection{RMSProp}

% =========================================================
% 1. EXPLICACIÓN GRÁFICA
% =========================================================
\begin{frame}{RMSProp: Explicación Gráfica}

\footnotesize

\begin{columns}[T,onlytextwidth]

\column{0.52\textwidth}

\begin{center}
\begin{tikzpicture}[scale=0.58]

    \draw[->, GrisTexto] (-2.6,0) -- (2.8,0) node[right] {$x$};
    \draw[->, GrisTexto] (0,-2.0) -- (0,2.2) node[above] {$y$};

    \draw[CelesteBrillante, thick] (0,0) ellipse (2.1 and 1.25);
    \draw[CelesteBrillante, thick] (0,0) ellipse (1.55 and 0.92);
    \draw[CelesteBrillante, thick] (0,0) ellipse (1.0 and 0.58);
    \draw[CelesteBrillante, thick] (0,0) ellipse (0.45 and 0.26);

    \fill[NaranjaBase] (2.0,1.35) circle (1.7pt);
    \fill[NaranjaBase] (1.35,0.72) circle (1.7pt);
    \fill[NaranjaBase] (0.88,0.36) circle (1.7pt);
    \fill[NaranjaBase] (0.42,0.14) circle (1.7pt);
    \fill[VerdeNeon] (0.10,0.03) circle (2pt);

    \draw[NaranjaBase, thick, ->] (2.0,1.35) -- (1.35,0.72);
    \draw[NaranjaBase, thick, ->] (1.35,0.72) -- (0.88,0.36);
    \draw[NaranjaBase, thick, ->] (0.88,0.36) -- (0.42,0.14);
    \draw[NaranjaBase, thick, ->] (0.42,0.14) -- (0.10,0.03);

\end{tikzpicture}

\textcolor{VerdeNeon}{Trayectoria hacia el mínimo}
\end{center}

\column{0.45\textwidth}

\textbf{¿Qué está ocurriendo?}

\begin{itemize}
    \item RMSProp guarda un promedio reciente de $g^2$.
    \item Gradiente grande $\Rightarrow$ paso menor.
    \item Gradiente pequeño $\Rightarrow$ paso relativamente mayor.
    \item Los gradientes antiguos pierden peso.
\end{itemize}

\[
s_k=\gamma s_{k-1}+(1-\gamma)g_k^2
\]

\[
x_{k+1}
=
x_k-
\frac{\alpha g_k}{\sqrt{s_k}+\epsilon}
\]

\end{columns}

\end{frame}

% =========================================================
% 2. PSEUDOCÓDIGO
% =========================================================
\begin{frame}{RMSProp: Pseudocódigo}

\small

\begin{algorithmic}[1]

\Require $x,\alpha,\gamma,\epsilon$

\State $s\gets0$

\While{no se cumpla la condición de parada}

    \State $g\gets\nabla f(x)$

    \State
    $s\gets
    \gamma s+(1-\gamma)(g\odot g)$

    \State
    $x\gets
    x-\dfrac{\alpha g}{\sqrt{s}+\epsilon}$

\EndWhile

\State \Return $x$

\end{algorithmic}

\vspace{0.35cm}

\begin{block}{Interpretación}
Un gradiente grande aumenta $s$ y reduce el paso.
Un gradiente pequeño permite un movimiento relativamente mayor.
\end{block}

\begin{center}
\footnotesize
$\gamma=0.9$:
90\% historial + 10\% gradiente actual.
\end{center}

\end{frame}


% =========================================================
% 3. EJEMPLO
% =========================================================
\begin{frame}{RMSProp: Ejemplo Paso a Paso}

\small

\[
f(x)=\frac12x^2,
\qquad
f'(x)=x
\]

\[
x_0=2,\quad
\alpha=0.1,\quad
\gamma=0.9,\quad
\epsilon=0.01,\quad
s_0=0
\]

\vspace{0.15cm}

\begin{center}
\scriptsize
\begin{tabular}{c|c|c|c|c}
\toprule
$k$ & $x_k$ & $g_k$ & $s_{k+1}$ & $x_{k+1}$ \\
\midrule
0 & 2.0000 & 2.0000 & 0.4000 & 1.6887 \\
1 & 1.6887 & 1.6887 & 0.6452 & 1.4810 \\
2 & 1.4810 & 1.4810 & 0.8000 & 1.3173 \\
\bottomrule
\end{tabular}
\end{center}

\vspace{0.2cm}

\footnotesize
\textbf{Primera iteración:}

\[
g_0=2
\qquad
s_1=0.9(0)+0.1(2^2)=0.4
\]

\[
x_1
=
2-\frac{0.1(2)}
{\sqrt{0.4}+0.01}
\approx1.6887
\]

\begin{center}
\textcolor{VerdeNeon}{
$2\rightarrow1.6887\rightarrow1.4810\rightarrow1.3173$
}
\end{center}

\end{frame}


% =========================================================
% 4. CÓDIGO
% =========================================================
\begin{frame}[fragile]{RMSProp: Código}

\begin{lstlisting}[
language=Python,
basicstyle=\ttfamily\scriptsize,
showstringspaces=false
]
def rmsprop(grad, x, alpha=.1,
            gamma=.9, eps=.01, n=10):
    s = 0.0

    for _ in range(n):
        g = grad(x)
        s = gamma*s + (1-gamma)*g**2
        x -= alpha*g/(s**0.5 + eps)

    return x


def grad(x):
    return x


print(rmsprop(grad, 2.0))
\end{lstlisting}

\end{frame}